% Initial draft of multi-draft watermark experiment results.
% Data source: 0505 multi_draft_n1000_efpr_4cell_multiT (n=1000 per cell, L=4,
%   B in {2,4,6,8}, decoders = MPFR (mpfr_torchgen_cached) and Invariant
%   (invariant_multi), 4 cells = {Qwen2.5-7B, Vicuna-7b-v1.5} x {CNN, ELi5}).
% Detection: Aaronson Gamma-tail at FPR alpha=1%, default window T=64.
% TPR-vs-T results are presented in a separate figure section (\ref{sec:tpr-figure}).

\section{Experiments}
\label{sec::experiments}

We evaluate whether the proposed Poisson-process sampler improves the empirical
frontier between speculative decoding efficiency and watermark detectability, while
preserving generation quality. Unless otherwise stated, we use Qwen2.5-7B-Instruct as
the target model and Qwen2.5-0.5B-Instruct as the draft model on
\textsc{CNN/DailyMail}. Additional model--dataset pairs are reported in the appendix.

\paragraph{Baselines and evaluation protocol.}
We compare against four groups of baselines. First, \textsc{basic\_uwm} denotes
autoregressive Gumbel-style unbiased watermarking without speculation. It serves as
a watermark-only reference for the detection signal achievable when no speculative
verification is used. Second, we compare with the speculative watermarking baselines
of \citet{hu2024inevitable}: \textsc{VSpS}, \textsc{MWS}, and \textsc{MSE}.
Here \textsc{VSpS} applies vanilla speculative sampling together with the unbiased
watermarking rule; \textsc{MWS} is designed to maintain watermark strength, but
typically pays a speculative-acceptance cost; and \textsc{MSE} is designed to
maintain sampling efficiency, but typically weakens the watermark signal. These
methods represent different points on the previously reported
efficiency--watermark trade-off. Third, for fixed-FPR detection comparisons, we also
include \textsc{MSE\_PSEUDO}, the pseudo-randomized variant of
\citet{he2026improving}, which improves the trade-off by injecting additional
pseudorandomness into the speculative acceptance or residual-sampling mechanism.
Fourth, in the multi-draft setting, we compare \textsc{MPFR} against the
unwatermarked invariant multi-draft decoder of \citet{rowan2025list}, which is the
closest list-coupling baseline for speculative efficiency. We also include
no-watermark variants of our methods to separate the cost of watermarking from the
cost of the Poisson coupling itself.

We report three classes of metrics. For speculative efficiency, we use average
accepted tokens per speculative step (AATPS), and report token rate in the appendix
as a system-level diagnostic. For watermark detectability, we report ANLPPT-U and
TPR@1\%FPR. The TPR threshold is calibrated on non-watermarked generations to obtain
a 1\% false-positive rate and then evaluated on watermarked generations. For
generation quality, we report target-model log perplexity (LPPL), together with
task metrics such as ROUGE-L when references are available. This separation is
important: AATPS measures the algorithmic acceptance behavior, ANLPPT-U and
TPR@1\%FPR measure detector-visible watermark evidence, and LPPL audits whether the
target-model marginal is preserved.

\paragraph{Single-draft efficiency--watermark frontier.}
Figure~\ref{fig::single_draft_qwen_cnn_new} shows the single-draft
efficiency--watermark frontier. The prior speculative watermarking baselines exhibit
the expected trade-off: methods that preserve stronger watermark evidence tend to
sacrifice speculative efficiency, while methods optimized for efficiency produce
weaker watermark evidence. In contrast, \textsc{PFR} lies in the upper-right region
of the plot, achieving both high AATPS and strong ANLPPT-U. Moreover,
\textsc{PFR} and its no-watermark variant have nearly identical AATPS, indicating
that the watermark signal is introduced through the target-side keyed race rather
than through an acceptance-penalizing reweighting step. Additional single-draft
results on other model--dataset pairs are provided in Appendix~\ref{sec:single_draft}.

\paragraph{Multi-draft scaling.}
We next evaluate whether the same construction scales to multiple drafts. We compare
\textsc{MPFR} with the unwatermarked invariant multi-draft decoder of
\citet{rowan2025list} using lookahead $L=4$ and draft counts $B\in\{2,4,6,8\}$.
Table~\ref{tab:headline} reports the representative multi-draft setting on
AATPS, token rate (TR), and ANLPPT-U.  As single-draft anchors we report
\textsc{pfr} (our watermarked single-draft method) and \textsc{pfr\_no\_wm}
(the same algorithm with the watermark key removed); the latter is the natural
single-draft baseline to which the unwatermarked Invariant of \citet{rowan2025list}
is implicitly compared in their original paper.  \textsc{MPFR} matches Invariant
in AATPS across draft counts and pays only a small TR overhead from watermark
labelling, while its ANLPPT-U sits at the level of \textsc{pfr} and far above the
two no-watermark controls (Invariant and \textsc{pfr\_no\_wm}).  Full results for
other model--dataset pairs, the additional speculative-watermarking baselines
\textsc{MWS}/\textsc{MSE} of \citet{hu2024inevitable}, the remaining ANLPPT
variants (Li, PL), and LPPL are deferred to Appendix~\ref{sec:multi_draft}.

\begin{table}[t]
\centering
\small
\caption{
Multi-draft scaling and single-draft anchors on Qwen2.5-7B/0.5B,
CNN/DailyMail, $n=1000$.  \textsc{MPFR} preserves AATPS relative to the
unwatermarked Invariant decoder of \citet{rowan2025list}, pays a
$\sim\!5\%$ TR overhead from watermark labelling, and maintains
ANLPPT-U at the level of the watermarked single-draft \textsc{pfr}.
\textsc{pfr\_no\_wm} is the same single-draft algorithm with the watermark
key removed; it is the natural single-draft baseline against which
Invariant is implicitly compared in \citet{rowan2025list}.
}
\label{tab:headline}
\begin{tabular}{lcccc}
\toprule
\multirow{2}{*}{Method}
& \multicolumn{4}{c}{AATPS $\uparrow$} \\
\cmidrule(lr){2-5}
& $B=2$ & $B=4$ & $B=6$ & $B=8$ \\
\midrule
\textsc{MPFR} (ours, watermarked)  & 2.80 & 3.12 & 3.30 & 3.41 \\
Invariant multi-draft (no wm)      & 2.80 & 3.11 & 3.28 & 3.39 \\
\midrule
\multicolumn{5}{l}{\textit{Token rate (tok/s) $\uparrow$}} \\
\midrule
\textsc{MPFR} (ours, watermarked)  & 23.18 & 24.89 & 25.39 & 25.17 \\
Invariant multi-draft (no wm)      & 25.03 & 26.63 & 26.79 & 26.38 \\
\midrule
\multicolumn{5}{l}{\textit{Watermark detectability: ANLPPT-U $\uparrow$}} \\
\midrule
\textsc{MPFR} (ours, watermarked)  & 0.052 & 0.052 & 0.052 & 0.053 \\
Invariant multi-draft (no wm)      & 0.003 & 0.002 & 0.002 & 0.002 \\
\midrule
\multicolumn{5}{l}{\textit{Single-draft anchors}} \\
\midrule
\textsc{pfr} (watermarked)         & \multicolumn{4}{c}{AATPS: 2.45 \quad TR: 20.76 \quad ANLPPT-U: 0.050} \\
\textsc{pfr\_no\_wm} (no wm)       & \multicolumn{4}{c}{AATPS: 2.44 \quad TR: 21.56 \quad ANLPPT-U: 0.002} \\
\bottomrule
\end{tabular}
\end{table}

\paragraph{Detection at fixed false-positive rate.}
ANLPPT-U measures average evidence per token, while TPR@1\%FPR measures operational
detectability under a calibrated false-positive budget. Figure~\ref{fig:drafter_invariance}
(left) compares detection power across methods, including the pseudo-randomized
\textsc{MSE\_PSEUDO} variant of \citet{he2026improving}. \textsc{PFR} and
\textsc{MPFR} closely track the autoregressive watermark reference and the
strong-watermark baseline \textsc{MWS}, while efficiency-oriented baselines such as
\textsc{MSE} produce weaker detection. \textsc{MSE\_PSEUDO} improves over
\textsc{MSE}, but still remains below the Poisson-race methods in our evaluated
settings. This supports the claim that the Poisson construction preserves
detector-visible watermark evidence while still benefiting from speculative
verification.

\begin{figure}[t]
\centering
\includegraphics[width=\linewidth]{images/drafter_invariance_composite.png}
\caption{
\textbf{Watermark detection and drafter substitution.}
\textbf{Left:} TPR@1\%FPR on the default drafter. \textsc{PFR}/\textsc{MPFR}
achieve detection power close to the autoregressive watermark reference
\textsc{basic\_uwm} and the strong-watermark baseline \textsc{MWS}. The
pseudo-randomized \textsc{MSE\_PSEUDO} baseline improves over \textsc{MSE} but
does not close the gap to the Poisson-race methods.
\textbf{Right:} TPR@1\%FPR under four drafter conditions: the default drafter
$D_0$, a model swap to Qwen2.5-1.5B-Instruct $D_1$, a sharper drafter
$D_2$ with $T_{\mathrm{drafter}}=0.5$, and a more diffuse drafter $D_3$ with
$T_{\mathrm{drafter}}=1.5$. Invariant methods produce nearly overlapping curves,
whereas drafter-sensitive methods show larger drift.
}
\label{fig:drafter_invariance}
\end{figure}

\paragraph{Drafter-substitution ablation.}
Finally, we test whether the watermark signal remains stable when the proposal
mechanism changes. We keep the target model, watermark key, and evaluation prompts
fixed, and vary only the drafter by either changing its model scale or its sampling
temperature. Figure~\ref{fig:drafter_invariance} (right) shows that \textsc{PFR}
has at most one percentage point of TPR drift at
$T_{\mathrm{eval}}\in\{64,128\}$ across the four drafter conditions. In contrast,
\textsc{MSE} drifts by $12$--$22$ percentage points, and \textsc{MSE\_PSEUDO}
drifts by $6$--$11$ percentage points across drafter conditions. This behavior is
consistent with the intended role of target-side keyed Poisson races: the drafter can
affect how many tokens are verified in a speculative block, but the watermark
statistic of the emitted tokens is tied to the realized prefix and the target-side
keyed randomness. We further verify that generation quality is stable across
decoders; full LPPL and ROUGE tables are reported in Appendix~\ref{app:ablation}.



\appendix
\subsubsection{Multi-draft experiments: full per-cell results}
\label{sec:multi_draft}

\paragraph{Setup.}
The main-text multi-draft evaluation
(Section~\ref{sec::experiments}, Table~\ref{tab:headline}) reports a
single representative cell.  Here we extend the evaluation to the
$2{\times}2$ matrix of \{Qwen2.5-7B-Instruct, Vicuna-7b-v1.5\}
$\times$ \{CNN/DailyMail, ELi5\}, using target/drafter pairs
\texttt{Qwen2.5-7B-Instruct} / \texttt{Qwen2.5-0.5B-Instruct} and
\texttt{lmsys/vicuna-7b-v1.5} / \texttt{double7/vicuna-68m}.  Each cell
uses $n{=}1000$ prompts, lookahead $L{=}4$, draft counts
$B\in\{2,4,6,8\}$, and $\mathrm{top\_k}=50,\ \mathrm{top\_p}=1.0,\
T_{\mathrm{eval}}=1.0$.

\paragraph{Decoders.}
\textsc{MPFR} denotes our keyed multi-draft Poisson sampler;
\textsc{Invariant} denotes the unwatermarked
\emph{InvariantMultiDraftStrategy} of \citet{rowan2025list} run in the
same harness on the same prompt set.  The two decoders share the
draft/verify mechanism and differ only in whether the per-context
watermark key is applied.

\paragraph{Per-cell tables.}
Tables~\ref{tab:full-qwen-cnn}--\ref{tab:full-vicuna-eli5} report all
quality and watermark-strength metrics per (decoder, $B$).  Columns:
average accepted tokens per step (AATPS); token rate in tokens/sec
(TR); ANLPPT-\{U,\,Li,\,PL\} (per-token recovery scores under the
three score variants); LPPL (per-token log-perplexity of the generated
tokens under the target model).  TPR-vs-$T$ curves are reported as a
figure in Section~\ref{sec:tpr-figure}.

\begin{table}[h]
\centering
\small
\caption{Qwen2.5-7B-Instruct $\times$ CNN/DailyMail.}
\label{tab:full-qwen-cnn}
\begin{tabular}{l c c c c c c c}
\toprule
Decoder & $B$ & AATPS & TR & U & Li & PL & LPPL \\
\midrule
MPFR      & 2 & 2.799 & 23.18 & 0.052 & 0.037 & 0.055 & 0.475 \\
MPFR      & 4 & 3.124 & 24.89 & 0.052 & 0.037 & 0.055 & 0.474 \\
MPFR      & 6 & 3.298 & 25.39 & 0.052 & 0.037 & 0.055 & 0.476 \\
MPFR      & 8 & 3.410 & 25.17 & 0.053 & 0.037 & 0.055 & 0.476 \\
Invariant & 2 & 2.804 & 25.03 & 0.003 & 0.003 & 0.002 & 0.487 \\
Invariant & 4 & 3.114 & 26.63 & 0.002 & 0.002 & 0.002 & 0.481 \\
Invariant & 6 & 3.275 & 26.79 & 0.002 & 0.003 & 0.002 & 0.482 \\
Invariant & 8 & 3.387 & 26.38 & 0.002 & 0.003 & 0.002 & 0.479 \\
\bottomrule
\end{tabular}
\end{table}

\begin{table}[h]
\centering
\small
\caption{Qwen2.5-7B-Instruct $\times$ ELi5.}
\label{tab:full-qwen-eli5}
\begin{tabular}{l c c c c c c c}
\toprule
Decoder & $B$ & AATPS & TR & U & Li & PL & LPPL \\
\midrule
MPFR      & 2 & 2.563 & 22.38 & 0.104 & 0.076 & 0.108 & 0.704 \\
MPFR      & 4 & 2.919 & 24.48 & 0.104 & 0.076 & 0.107 & 0.703 \\
MPFR      & 6 & 3.111 & 25.54 & 0.105 & 0.076 & 0.108 & 0.706 \\
MPFR      & 8 & 3.237 & 25.83 & 0.105 & 0.076 & 0.108 & 0.704 \\
Invariant & 2 & 2.519 & 23.18 & 0.002 & 0.002 & 0.002 & 0.750 \\
Invariant & 4 & 2.852 & 25.75 & 0.002 & 0.002 & 0.002 & 0.733 \\
Invariant & 6 & 3.043 & 27.11 & 0.002 & 0.002 & 0.002 & 0.730 \\
Invariant & 8 & 3.158 & 27.96 & 0.002 & 0.002 & 0.002 & 0.725 \\
\bottomrule
\end{tabular}
\end{table}

\begin{table}[h]
\centering
\small
\caption{Vicuna-7b-v1.5 $\times$ CNN/DailyMail.}
\label{tab:full-vicuna-cnn}
\begin{tabular}{l c c c c c c c}
\toprule
Decoder & $B$ & AATPS & TR & U & Li & PL & LPPL \\
\midrule
MPFR      & 2 & 2.422 & 52.10 & 0.022 & 0.017 & 0.022 & 0.270 \\
MPFR      & 4 & 2.664 & 48.28 & 0.022 & 0.017 & 0.023 & 0.271 \\
MPFR      & 6 & 2.803 & 44.44 & 0.022 & 0.017 & 0.023 & 0.270 \\
MPFR      & 8 & 2.899 & 40.76 & 0.022 & 0.017 & 0.023 & 0.270 \\
Invariant & 2 & 2.478 & 58.15 & 0.002 & 0.002 & 0.002 & 0.287 \\
Invariant & 4 & 2.690 & 53.12 & 0.002 & 0.002 & 0.002 & 0.286 \\
Invariant & 6 & 2.831 & 48.41 & 0.002 & 0.003 & 0.002 & 0.283 \\
Invariant & 8 & 2.925 & 44.01 & 0.002 & 0.002 & 0.002 & 0.288 \\
\bottomrule
\end{tabular}
\end{table}

\begin{table}[h]
\centering
\small
\caption{Vicuna-7b-v1.5 $\times$ ELi5.}
\label{tab:full-vicuna-eli5}
\begin{tabular}{l c c c c c c c}
\toprule
Decoder & $B$ & AATPS & TR & U & Li & PL & LPPL \\
\midrule
MPFR      & 2 & 1.889 & 45.89 & 0.058 & 0.043 & 0.061 & 0.504 \\
MPFR      & 4 & 2.133 & 47.43 & 0.058 & 0.042 & 0.061 & 0.504 \\
MPFR      & 6 & 2.278 & 47.42 & 0.058 & 0.042 & 0.062 & 0.503 \\
MPFR      & 8 & 2.377 & 46.42 & 0.059 & 0.043 & 0.062 & 0.504 \\
Invariant & 2 & 1.884 & 50.29 & 0.002 & 0.002 & 0.002 & 0.555 \\
Invariant & 4 & 2.108 & 53.42 & 0.002 & 0.002 & 0.002 & 0.552 \\
Invariant & 6 & 2.248 & 55.08 & 0.002 & 0.003 & 0.002 & 0.546 \\
Invariant & 8 & 2.361 & 56.00 & 0.003 & 0.003 & 0.002 & 0.539 \\
\bottomrule
\end{tabular}
\end{table}


\paragraph{Speculative-acceptance parity.}
Within each cell, \textsc{MPFR}'s AATPS matches \textsc{Invariant}'s
within $\pm 1\%$ at every $B \in \{2,4,6,8\}$, indicating that the
target-side keyed-race construction adds no per-step acceptance
penalty.  Realized token rate (TR) tracks AATPS up to a fixed
Python-side label-hashing overhead whose relative cost depends on the
target/drafter compute ratio: with the small \texttt{vicuna-68m}
drafter the overhead is a larger fraction of per-step time, and TR is
$\sim\!10\%$ below \textsc{Invariant}, while on the
\texttt{Qwen2.5-7B/0.5B} pair the gap narrows to $\sim\!5\%$.

\paragraph{Stability of the watermark signal across $B$.}
$\mathrm{ANLPPT}_U$, $\mathrm{ANLPPT}_{Li}$, and $\mathrm{ANLPPT}_{PL}$
for \textsc{MPFR} remain within $\pm 0.001$ across $B\in\{2,4,6,8\}$ on
every cell of Tables~\ref{tab:full-qwen-cnn}--\ref{tab:full-vicuna-eli5}.
The per-token detection signal therefore does not dilute with the
number of drafts, in contrast to naive multi-draft watermarking, which
would re-introduce the efficiency--detectability trade-off of
\citet{hu2024inevitable} at the multi-draft level.

\paragraph{Distortion-free audit.}
LPPL of \textsc{MPFR} outputs falls within $\pm 0.013$ nats/token of
\textsc{Invariant} on every cell.  This is consistent with per-token
marginal preservation and confirms that the watermark does not
introduce a measurable shift in target-model log-perplexity in the
multi-draft regime.